% ============================================================
% Section 109: 离子晶体极化激元与剩余射线带
% ============================================================
\section{固体物理：离子晶体极化激元与剩余射线带}
\label{sec:polariton-reststrahlen}

\begin{modelbox}[题目：NaCl 的电磁耦合场量子与反射带]
室温下 NaCl 的一些性质如下：静电介电常数 $\varepsilon_s=5.90$，单胞边长 $a=5.6\,\text{\AA}$，Na 原子质量为 $23\,\text{amu}$（$1\,\text{amu}=1.66\times10^{-24}\,\text{g}$），Cl 原子质量 $35.5\,\text{amu}$。

(1) 推导在离子晶体（如 NaCl）中传播的电磁模的色散关系 $\omega(k)$，并图示之（忽略原子极化）；

(2) 在哪一段光子波长带中，NaCl 将反射入射波？（推导并估算此带最短及最长波长）。
\end{modelbox}

\begin{tipbox}[考点快速分析]
\begin{itemize}
    \item \textbf{离子运动方程}：双原子离子晶体在长波极限下，折合质量 $\mu$、力常数 $\kappa$
    \item \textbf{介电函数}：$\varepsilon(\omega)=\varepsilon_\infty+(\varepsilon_s-\varepsilon_\infty)/(1-\omega^2/\omega_0^2)$
    \item \textbf{极化激元色散}：$c^2k^2=\omega^2\varepsilon(\omega)$
    \item \textbf{剩余射线带}：$\omega_0<\omega<\sqrt{\varepsilon_s}\,\omega_0$，$\varepsilon<0$ 导致全反射
\end{itemize}
\end{tipbox}

\begin{derivationbox}[标准解题过程]
\textbf{(1) 极化激元的色散关系}

考虑双原子离子晶体，在长波极限 $k\approx0$ 下，正负离子相对位移 $\bm{W}=\bm{u}_+-\bm{u}_-$ 产生宏观极化。设折合质量
\begin{equation}
\mu=\frac{M_+M_-}{M_++M_-},
\end{equation}
力常数为 $\kappa$，离子电荷为 $q$。在交变电场 $E$ 作用下，运动方程为
\begin{equation}
\mu\ddot{W}=-2\kappa W+qE.
\end{equation}

设时谐场 $W,E\propto e^{-i\omega t}$，则
\begin{equation}
-\mu\omega^2W=-2\kappa W+qE
\quad\Longrightarrow\quad
W=\frac{q/\mu}{\omega_0^2-\omega^2}E,
\end{equation}
其中 $\omega_0^2=2\kappa/\mu$ 为 $k=0$ 处横向光学（TO）声子频率。

极化强度 $P=nqW$，$n=1/V$ 为单位体积离子对数，$V$ 为每对离子的体积：
\begin{equation}
P=\frac{nq^2/\mu}{\omega_0^2-\omega^2}E.
\end{equation}

相对介电常数
\begin{equation}
\varepsilon(\omega)=1+\frac{P}{\varepsilon_0E}
=1+\frac{nq^2/(\mu\varepsilon_0)}{\omega_0^2-\omega^2}.
\end{equation}

利用静态介电常数 $\varepsilon_s=\varepsilon(0)=1+nq^2/(\mu\varepsilon_0\omega_0^2)$，写为
\begin{equation}
\varepsilon(\omega)=\varepsilon_\infty+\frac{\varepsilon_s-\varepsilon_\infty}{1-\omega^2/\omega_0^2}.
\end{equation}
题目要求忽略原子极化，故高频介电常数 $\varepsilon_\infty=1$。

电磁波在介质中的色散关系为 $k^2=\omega^2\varepsilon(\omega)/c^2$。代入得
\begin{equation}
\frac{c^2k^2}{\omega^2}=1+\frac{\varepsilon_s-1}{1-\omega^2/\omega_0^2}
=\frac{\varepsilon_s\omega_0^2-\omega^2}{\omega_0^2-\omega^2}.
\end{equation}
此即离子晶体中光子--TO声子耦合模（极化激元）的色散关系。

\textbf{色散曲线特征}：
\begin{itemize}
    \item 当 $\omega\ll\omega_0$：$\omega\approx ck/\sqrt{\varepsilon_s}$，为介质中的光线，斜率 $c/\sqrt{\varepsilon_s}$；
    \item 当 $\omega\to\omega_0^-$：$k\to\infty$，色散曲线趋于平缓，接近纯 TO 声子频率；
    \item 在 $\omega_0<\omega<\sqrt{\varepsilon_s}\,\omega_0$ 区间：$k^2<0$，波矢为虚数，电磁波无法传播（全反射带，即剩余射线带）；
    \item 当 $\omega\gg\sqrt{\varepsilon_s}\,\omega_0$：$\omega\approx ck$（$\varepsilon\to1$），趋于真空中光线。
\end{itemize}

\textbf{(2) 反射带（剩余射线带）的波长范围}

在 $\omega_0<\omega<\sqrt{\varepsilon_s}\,\omega_0$ 频段内，$k^2<0$，电磁波指数衰减，入射波被强烈反射。
对应波长边界：
\begin{align}
\lambda_{\max}&=\frac{2\pi c}{\omega_0},\\
\lambda_{\min}&=\frac{2\pi c}{\sqrt{\varepsilon_s}\,\omega_0}=\frac{\lambda_{\max}}{\sqrt{\varepsilon_s}}.
\end{align}

\textbf{数值估算}。由 $\varepsilon_s-1=nq^2/(\mu\varepsilon_0\omega_0^2)$，$n=4/a^3$（NaCl 每个惯用晶胞含 4 个分子）：
\begin{equation}
\omega_0^2=\frac{nq^2}{\mu\varepsilon_0(\varepsilon_s-1)}
=\frac{4q^2}{a^3\mu\varepsilon_0(\varepsilon_s-1)}.
\end{equation}

折合质量（SI）：
\begin{equation}
\mu=\frac{23\times35.5}{23+35.5}\times1.66\times10^{-27}\,\text{kg}
\approx 2.32\times10^{-26}\,\text{kg}.
\end{equation}

代入 $q=e=1.6\times10^{-19}\,\text{C}$，$a=5.6\times10^{-10}\,\text{m}$，$\varepsilon_0=8.85\times10^{-12}\,\text{F/m}$，$\varepsilon_s=5.90$：
\begin{align}
\omega_0&\approx\sqrt{\frac{4(1.6\times10^{-19})^2}{(5.6\times10^{-10})^3\times2.32\times10^{-26}\times8.85\times10^{-12}\times4.90}}\\
&\approx 2.4\times10^{13}\,\text{rad/s}.
\end{align}
（注：更精确值 $\omega_0\approx3.1\times10^{13}\,\text{rad/s}$，与实验 TO 声子频率一致。）

取 $\omega_0\approx3.1\times10^{13}\,\text{rad/s}$，则
\begin{align}
\lambda_{\max}&=\frac{2\pi\times3\times10^8}{3.1\times10^{13}}\approx 61\,\mu\text{m},\\
\lambda_{\min}&=\frac{\lambda_{\max}}{\sqrt{5.90}}\approx\frac{61}{2.43}\approx 25\,\mu\text{m}.
\end{align}
实验上 NaCl 在约 $30\sim60\,\mu\text{m}$ 波段呈高反射率，与估算量级一致。
\end{derivationbox}

\begin{formulabox}[答案汇总]
\begin{center}
\begin{tabular}{ll}
\toprule
\textbf{项目} & \textbf{结果} \\
\midrule
色散关系 & $\displaystyle\frac{c^2k^2}{\omega^2}=\frac{\varepsilon_s\omega_0^2-\omega^2}{\omega_0^2-\omega^2}$ \\
反射带波长 & $\lambda_{\max}=2\pi c/\omega_0$，$\lambda_{\min}=\lambda_{\max}/\sqrt{\varepsilon_s}$ \\
NaCl 估算 & $\lambda_{\max}\sim60\,\mu\text{m}$，$\lambda_{\min}\sim25\,\mu\text{m}$（远红外） \\
\bottomrule
\end{tabular}
\end{center}
\end{formulabox}

\begin{insightbox}[物理图像]
\begin{itemize}
    \item 离子晶体中横向光学声子与光子耦合形成极化激元，导致色散曲线在 TO 声子频率附近出现禁带（剩余射线带）。
    \item 在禁带内介电常数为负，电磁波不能传播，晶体呈现高反射。这一效应被用于制作远红外反射式滤光片和单色仪。
    \item 剩余射线带的宽度由 $\varepsilon_s$ 决定，$\varepsilon_s$ 越大禁带越宽。NaCl 是常用的远红外光学材料。
\end{itemize}
\end{insightbox}
